\section{One conditioning step: arm MU}\label{sec:arm-mu}

Arm MU hands the proposer one parent and one instruction, with no iteration. Given a
parent at the family argmax the proposer keeps it, 26 of 26, and none of those outputs sits
on the anchor. Given its own anchor as parent it leaves, 17 of 28 moving above it.

\textbf{Design.} Each call is the bare prompt, an existing packing, and the instruction to
increase the sum of radii.
The \emph{anchor parent} is built at the model's own unconditioned anchor value $T(k^*, N)$.
The \emph{argmax parent} is built at the family argmax $V(k^*{-}1, m)$, which at these cells
is provably better than the anchor. The \emph{off-family parent} is a diagonal-chain packing
scoring far below either. The registration and the supplement name them A\_anchor, B\_rival
and C\_offfamily. The cells are $N = 13$, 21 and 31. Three parent conditions in three cells,
at $n = 15$ samples each, give 135 invocations.

\textbf{Registered predictions.} P-MU1 predicted an on-prediction rate under the anchor parent
of at least 50\%. P-MU2 predicted an on-prediction rate under the argmax parent of at least
35\%. P-MU3 predicted $k$-inheritance under the argmax parent below 50\%. P-MU4 predicted that
under the off-family parent the modal valid output sits on the $k^*$-grid, scored per cell.
The registered decision rule calls anchoring dissolved when the pooled on-prediction rate
under the argmax parent is at most 20\% and the keep-or-improve rate is at least 50\%. The
falsifier fires on the same condition.

\textbf{The rule's two conditions under the argmax parent.} Both are met, and they are not
independent. The pooled on-prediction rate is 0 of 26 valid outputs. That meets
the rule's 20\% ceiling and fails P-MU2's 35\% floor. Keep-or-improve is 26 of 26, above
its floor of 50\%. The argmax parent outscores the anchor
$T(k^*, N)$ by 0.151 at $N = 13$. The gap is 0.159 at $N = 21$. It is 0.165 at $N = 31$. Each
gap is more than seventy times the $2\times10^{-3}$ matching window, so an output that keeps or
improves on that parent cannot sit on the anchor. The 26 of 26 entails the 0 of 26,
and the rule reduces to one condition.

Keep-or-improve is almost entirely keep. Of the 26, 25 return the parent's score within
$10^{-6}$, and one exceeds it by $1.2\times10^{-6}$. In
coordinates, 18 reproduce the parent to within $10^{-6}$ on every coordinate. The
$1.2\times10^{-6}$ excess is among them, from radii perturbed below that tolerance. One
further output is a mirror image of the parent. The other 7 relocate
filler circles while keeping the radius multiset and the parent's sum to seven decimals. All
26 inherit the parent's grid order too, so P-MU3 is disconfirmed.

The independent evidence is the anchor parent. Its on-prediction rate is 7 of 28,
below P-MU1's 50\% threshold, so the anchor is not self-stable as its own parent. Here
the model departs: 17 of 28 move above the anchor and 4 fall below it. Grid order is kept by
22 of 28, 79\% (60--90\%). The 6 that change it all move above the anchor, and they include
3 that reach the family argmax value. The falsifier negates the persistence prediction
over this pair of conditions, and no other was registered.

\textbf{The unconditioned rate on the same instrument.} Arm L's generation 0
(\S\ref{sec:arm-l}) supplies the no-parent reference on the same prompt bytes, channel and
wrapper. It puts 4 of 9 valid outputs on the anchor at $N = 13$. It puts 7 of 8
there at $N = 31$. The anchor parent gives 4 of 12 and 2 of 11 at those cells. Arm
L did not run at $N = 21$, so no unconditioned row exists there.

\textbf{P-MU4, held in two cells of three.} The anchor's grid order does reassert itself in the one
condition registration predicted. Given the off-family diagonal parent, 20 of 34
valid samples abandon the parent and emit a $k^*$-grid construction. The prediction holds
at $N = 21$ and at $N = 31$, and not
at $N = 13$.

\begin{center}
\begin{small}
\setlength{\tabcolsep}{3.5pt}
\begin{tabular}{@{}p{0.18\linewidth}cp{0.26\linewidth}p{0.15\linewidth}p{0.21\linewidth}p{0.06\linewidth}@{}}
\toprule
Parent & Valid & On-prediction (95\% Wilson) & Registered threshold & Score against parent & $k$-inher. \\
\midrule
None (arm L generation 0, same instrument) & 17 & 4/9 at $N = 13$ (44\%, [19--73\%]); 7/8 at $N = 31$ (88\%, [53--98\%]); no row at $N = 21$ & unconditioned reference, no bar & --- & --- \\
Argmax parent (a packing at the family argmax) & 26 & 0/26 (0\%, [0--13\%]) & dissolved if $\le$ 20\% & 26/26 (100\%, [87--100\%]; bar $\ge$ 50\%): 25 equal within $10^{-6}$, 1 above by $1.2\times10^{-6}$ & 26/26 \\
Anchor parent (the model's own anchor) & 28 & 7/28 (25\%, [13--43\%]); per cell 4/12, 1/5 and 2/11 at $N = 13$, 21 and 31 & P-MU1 predicted $\ge$ 50\%, not met & 17/28 above the anchor, 4/28 below & 22/28 \\
Off-family parent (diagonal chain, sum $\approx$ 0.6) & 34 & 20/34 abandon parent for $k^*$-grid (59\%, [42--74\%]); per cell 1/7 [3--51\%], 7/13 [29--77\%] and 12/14 [60--96\%] at $N = 13$, 21 and 31 & P-MU4: $k^*$-grid reasserts, scored per cell & --- & --- \\
\bottomrule
\end{tabular}

{\footnotesize Bracketed ranges are 95\% Wilson intervals on the rate they follow.}
\end{small}
\end{center}


\textbf{Clearance of the family argmax (post hoc, disclosed).} Arm MU's rows were scored
against the anchor, so the clearance question was never asked at registration. Recomputing
the argmax from the closed form gives $1.776142375$ at $N = 13$. At $N = 21$ it is
$2.258883476$. At $N = 31$ it is $2.748528137$. The 135 invocations give 88 valid outputs at
the primary tolerance. At $10^{-9}$ just 63 of the 135 are valid. Of the 88, 27 return a
sum nominally above the argmax: 9 within the $5\times10^{-8}$ granularity of a recorded sum,
18 beyond it, and 61 at or below. Seventeen of the eighteen are argmax-parent outputs
returning their parent's score within $1.2\times10^{-6}$. The eighteenth is the anchor-parent output at
$N = 21$ that moved onto the argmax.

The argmax parent's seven-decimal coordinates sum to $3.25\times10^{-7}$ above the
closed form at $N = 13$ and $2.63\times10^{-7}$ above it at $N = 31$. Thirteen of the
eighteen return that parent sum unchanged; their excess is the parent's rounding. The other
five sit up to $1.53\times10^{-6}$ above the argmax. All eighteen fail the $10^{-9}$ test on
overlapping circles, not on containment. Their largest single overlaps run from
$3.8\times10^{-9}$ to $4.3\times10^{-7}$, median $4.4\times10^{-8}$, each smaller than its
row's excess. Rounded or perturbed coordinates carry the excess; the $10^{-6}$ tolerance
admits the overlaps they create (per-row table and the granularity split of the 88,
Supplement~S2). Of the 63 valid at $10^{-9}$, four sit above the argmax by at most
$2.5\times10^{-8}$, inside the recorded-sum granularity, and none clears.
